\documentclass[11pt]{report}
\usepackage[a4paper,width=150mm,top=25mm,bottom=25mm]{geometry}
\usepackage[spanish]{babel}
\usepackage[utf8]{inputenc}
\usepackage{fancyhdr}
\usepackage{xhfill}
\usepackage{nameref}
\pagestyle{fancy}
\usepackage{tikz}
\usetikzlibrary{babel}

\usepackage{blindtext}
\usepackage{graphicx}
\usepackage{tabularx}
\usepackage{listingsutf8}
\usepackage{color}
\usepackage{enumitem}

\definecolor{mygreen}{rgb}{0,0.6,0}
\definecolor{mygray}{rgb}{0.5,0.5,0.5}
\definecolor{mymauve}{rgb}{0.58,0,0.82}
\renewcommand{\lstlistingname}{Listado}
\lstset{
  backgroundcolor=\color{white},   % choose the background color; you must add \usepackage{color} or \usepackage{xcolor}; should come as last argument
  basicstyle=\footnotesize,        % the size of the fonts that are used for the code
  breakatwhitespace=false,         % sets if automatic breaks should only happen at whitespace
  breaklines=true,                 % sets automatic line breaking
  captionpos=b,                    % sets the caption-position to bottom
  commentstyle=\color{mygreen},    % comment style
  deletekeywords={...},            % if you want to delete keywords from the given language
  escapeinside={\%*}{*)},          % if you want to add LaTeX within your code
  firstnumber=1,                   % start line enumeration with line 1
  frame=single,	                   % adds a frame around the code
  inputpath=../servidor/src/sistop/,
  inputencoding=utf8/latin1,
  keepspaces=true,                 % keeps spaces in text, useful for keeping indentation of code (possibly needs columns=flexible)
  keywordstyle=\color{blue},       % keyword style
  language=Java,                 % the language of the code
  morekeywords={*,...},            % if you want to add more keywords to the set
  numbers=none,                    % where to put the line-numbers; possible values are (none, left, right)
  numbersep=5pt,                   % how far the line-numbers are from the code
  numberstyle=\tiny\color{mygray}, % the style that is used for the line-numbers
  rulecolor=\color{black},         % if not set, the frame-color may be changed on line-breaks within not-black text (e.g. comments (green here))
  showspaces=false,                % show spaces everywhere adding particular underscores; it overrides 'showstringspaces'
  showstringspaces=false,          % underline spaces within strings only
  showtabs=false,                  % show tabs within strings adding particular underscores
  stepnumber=2,                    % the step between two line-numbers. If it's 1, each line will be numbered
  stringstyle=\color{mymauve},     % string literal style
  tabsize=2,	                   % sets default tabsize to 2 spaces
  title=\lstname,                   % show the filename of files included with \lstinputlisting; also try caption instead of title
  xleftmargin=15pt,
  xrightmargin=15pt
}
% \usepackage{algorithmicx}
\usepackage{algorithm}
\usepackage{algpseudocode}
% \usepackage{fontawesome}

\fancyhead{}
\fancyhead[R]{Obligatorio}
\fancyhead[L]{Capítulo \thechapter}



\makeatletter
\renewcommand{\ALG@name}{Algoritmo}
\makeatother

\algrenewcommand\algorithmicprocedure{\textbf{procedimiento}}
\algrenewcommand\algorithmicif{\textbf{si}}
\algrenewcommand\algorithmicthen{\textbf{entonces}}
\algrenewcommand\algorithmicelse{\textbf{si no}}
\algrenewcommand\algorithmicend{\textbf{fin}}
\algrenewcommand\algorithmicrequire{\textbf{Pre:}}
\algrenewcommand\algorithmicensure{\textbf{Post:}}
\algrenewcommand\algorithmicforall{\textbf{para cada}}
\algrenewcommand\algorithmicdo{\textbf{hacer}}
\algrenewcommand{\algorithmiccomment}[1]{\hfill$\langle$ #1}

\graphicspath{ {img/} }

\newcommand{\threadFig}[0]{
    \begin{tikzpicture}
      \vspace{5pt}
      \draw[very thick,color=black] plot [smooth,tension=1.5] coordinates{(0.3,0.5) (0.25,0.37) (0.3,0.25) (0.25,0.12)};
    \end{tikzpicture}
}

\newcommand{\BigO}[1]{\ensuremath{\mathrm{O}\bigl(#1\bigr)}}
\newcommand{\type}[1]{\mbox{\textbf{#1}}}
\newcommand{\call}[1]{\mbox{\texttt{#1}}}

\overfullrule=2cm

\begin{document}

\begin{titlepage}
    \begin{center}
        \includegraphics[width=0.7\textwidth]{ucu}

	\vspace*{5cm}

        \Huge
        \textbf{Obligatorio}

        \vspace{0.5cm}
        \LARGE
	Teoría de la Computación

        \vspace{3.5cm}

        Lois, Lucas

        Maresca, Santiago

        Silveira, Marco

        \vfill

        \Large
        Universidad Católica del Uruguay\\
        22 de Noviembre, 2019

    \end{center}
\end{titlepage}


\tableofcontents

\chapter{Recursos y Procesos Involucrados}
\section{Recursos}

\subsubsection{Asientos}

Cada asiento tiene una ubicación específica y única en el estadio.

\subsubsection{Lugares en cancha}

El tamaño de la cancha determina la cantidad de lugares, pero estos no tienen
cada uno una ubicación específica como los asientos.

\subsubsection{Esquema del estadio}

El esquema es la distribución física de los espacios (asientos y cancha) en él.

\subsubsection{Entradas}

Las entradas, asignadas a un espacio y un espectáculo en particular, poseen
además un precio determinado.

\subsubsection{Espectáculos}\label{recurso:espectaculo}

Los espectáculos tienen una fecha y hora, duración, nombre, así como entradas asociadas.

\subsubsection{Agencias, Organizadores de eventos y puntos de venta}

Estas entidades poseen la capacidad de funcionar como intermediarios en la venta
de entradas a usuarios finales. Deben ser capaces de identificarse ante la
plataforma como tales, para asegurar permisos diferentes.

En varios casos, según se dictamine al agregar el espectáculo a la plataforma,
estos podrían poseer preferencia o derecho sobre las entradas o un subconjunto
de estas.

\section{Procesos}

\subsubsection{Agregar un Espectáculo}

Al agregar un espectáculo al sistema debe darse información sobre el mismo, como
su nombre, categoría y duración.  Además, deben especificarse categorías y
precios de las entradas, según localidad.

\subsubsection{Obtener entradas disponibles}

Debe ser posible obtener un listado de los asientos o entradas disponibles para
un evento determinado.

\subsubsection{Identificación/Validación de usuarios}

Los usuarios del sistema (agencias, puntos de venta y usuarios de la aplicación o sitio web)
deben ser capaces de validar su identidad para asegurar un
funcionamiento seguro del mismo.

\subsubsection{Registro en la plataforma}

Los usuarios del sitio web o la aplicación deben ser capaces de registrarse
para poder luego realizar compras.

\subsubsection{Venta de entradas a privados}

Un usuario debe ser capaz de realizar la compra de una o más
entradas para un espectáculo determinado. Pueden aplicarse restricciones
(e.g.\ asientos para público con capacidades diferentes o asientos no-contiguos).

\subsubsection{Venta de entradas a mayoristas}

Como tales, y eventualmente solo si el espectáculo en particular lo permitiera, los mayoristas\footnote{Agencias u organizadores} serán
capaces de realizar la compra de un lote (cantidad definida por espectáculo) de entradas,
a un precio por entrada menor que en la venta a privados.

\subsubsection{Correo de confirmación al usuario al comprar entrada}

Cada cliente que reserva un asiento debe recibir un correo que confirme la compra que acaba de realizar. 
El correo debe contener la información del evento y la información de la entrada.
Esta sección no se implementará hasta la entrega del producto final con el fin de simplificar las simplificar
las simulaciones de distintos escenarios.
\
\chapter{Ordenamiento Justificado de los Criterios}\label{cap:criterios}


\section{Maximizar el ingreso total por entradas vendidas}
Sin importar el tipo de entradas (e.g.\ sector VIP o Golden Tickets, paquetes
Hospitality, etc.) la prioridad principal es la de obtener ingresos, maximizando
el monto total de ventas.

\section{Maximizar la ocupación del estadio.}
Para atraer espectáculos (garantizando su éxito en cuanto a espectadores) se
prioriza maximizar la ocupación del estadio, medida por la cantidad total de entradas vendidas,
independiente de su precio.

\section{Minimizar el tiempo de espera del usuario}

Con rápidos tiempos de respuesta y una experiencia sin inconvenientes el usuario
se siente satisfecho en su desempeño y es probable que vuelva a operar en el
sistema.


\section{Minimizar la cantidad de asientos únicos aislados libres.}\label{asientos-aislados}

En cada venta sería ideal que no resten asientos sin vecinos libres ya que su venta será difícil.

\begin{figure}
\centering{
\input{img/asientos.pdf_tex}
\caption{En \textcircled{} un asiento aislado. \(\times\) muestra asientos ocupados.}
}
\end{figure}

\section{Maximizar la satisfacción de personas discapacitadas.}
Dando acceso preferencial o reservando asientos.



\chapter{Problemas similares en Sistemas Operativos e información teórica}
\section{Procesos}

``Informalmente, un proceso es un programa en ejecución. Además de código, un proceso incluye también la actividad actual, que
está formada por el valor del contador del programa y el contenido de los registros del procesador.
Cada proceso se representa en el sistema operativo mediante su propio bloque de control de proceso.''\cite{silberschatz}.

Generalmente, un proceso también incluye una pila, que contiene datos temporales, y una sección de datos, 
que contiene las variables globales.
\cite{silberschatz}.

\subsection{Estados de un proceso}

Los procesos de un sistema pueden encontrarse (en principio) en uno de 5 estados, que reflejan su actividad actual.

\begin{description}
	\item[Nuevo:] El proceso está siendo creado en el sistema.
	\item[Listo:] El proceso está a la espera de que lo asignen a un procesador.
	\item[Ejecución:] Se están ejecutando sus instrucciones. Tiene control de la CPU.
	\item[Bloqueado:] El proceso está en espera de que se produzca un suceso.
	\item[Terminado:] El proceso ha terminado su ejecución.
\end{description}

\subsection{Planificación de Procesos}

En un sistema de multiprocesamiento, los procesos (y eventualmente los hilos de kernel) compiten por el control de la CPU.
Cuando la cantidad de tareas\footnote{Procesos o hilos de kernel} supera a la cantidad de procesadores es evidente que algunas
deben esperar a que se libere alguno de ellos. Entonces, cuando no se esta ejecutando, se encuentra en alguna cola de espera.
Con la multiprogramación se busca utilizar al sistema de manera eficiente, permitiendo la ejecución de múltiples tareas de la mejor forma posible.
A modo de ejemplo, cuando una tarea tiene que esperar por una operación de E/S se planifica a otra en el procesador.
El algoritmo que decide qué tarea obtendrá control de la CPU se dice un \emph{planificador a corto plazo}\cite{silberschatz}.

Una de los aspectos que se busca de un planificador a corto plazo es que el tiempo de ejecución sea mínimo, así como el de espera.
El tiempo de ejecución es el intervalo que va desde el instante en que se ordena la ejecución de una tarea hasta el instante en que se completa \cite{silberschatz}.
El de espera\label{tiempo-espera} se define como la suma de los períodos en que la tarea se encuentra en el estado \emph{Listo}.

\subsubsection{Esquemas de Planificación}

Con el fin de no tener una CPU desocupada el planificador asigna siempre una nueva tarea cuando la tarea que se encontraba en ejecución pasa al estado de \emph{Bloqueado} o \emph{Terminado} \cite{catedra-planificacion}.
Algunos planificadores deciden, además, elegir un nuevo proceso (o al menos permitirse hacerlo) en las transiciones:
\begin{itemize}
	\item de \emph{Bloqueado} a \emph{Listo}
	\item de \emph{Ejecución} a \emph{Listo} (e.g.\ Al ocurrir una interrupción)
\end{itemize}

Estos algoritmos se dicen algoritmos de esquema \emph{apropiativo}.
Aquellos que no sean apropiativos se dicen de esquema \emph{cooperativo} o \emph{sin desalojo}.


\subsubsection{First Come, First Serve}\label{fcfs}

El algoritmo \emph{FCFS} es el más simple de los planificadores a corto plazo.
Con este esquema se asigna control de la CPU a la primer tarea que lo solicite.
Su implementación consiste en tomar a la cola de trabajos preparados (aquellos en estado \emph{Listo}) como una cola \textbf{FIFO}.
Este algoritmo es uno de esquema \emph{cooperativo}.

A pesar de la simplicidad del algoritmo (o a costas de ella) la planificación FCFS posee dos cualidades no particularmente deseables.
En primer lugar, el tiempo medio de espera (ver página \pageref{tiempo-espera}) es bastante largo\cite{silberschatz}.
Además, presenta un comportamiento conocido con el nombre de \emph{efecto convoy}.

Los trabajos de un sistema informático pueden verse como una cadena de dos clases de estado: ráfaga de CPU y ráfaga de E/S.
En ciertas ocasiones puede suceder que muchas tareas (que finalizaron su ráfaga de E/S y están en espera de que se les asigne CPU)
se vean incapacitadas de continuar porque deben esperar a que \emph{un} proceso acabe libere la CPU.
Esto se conoce como efecto convoy, y da lugar a una utilización de la CPU y los
   dispositivos menor que si se permitiera a procesos mas cortos ejecutar primero.

\subsubsection{Planificación por prioridades}

En este algoritmo se asocia con cada tarea una prioridad, y al momento de elegir un nuevo proceso se toma a aquel de mayor prioridad.
El mismo puede ser de cualquier esquema, apropiativo o cooperativo.

La prioridad no está definida a-priori, y es algo que dependerá del sistema particular que esté bajo consideración.
Ella puede ser definida de dos maneras, interna o externamente.

La prioridad definida internamente utiliza alguna medida de la tarea para calcularla.
Un ejemplo de esto es la planificación SJF (\emph{Shortest Job First}) que estima
(típicamente mediante una media exponencial) la duración de la próxima ráfaga de CPU.
Esto le permite ejecutar el trabajo más corto primero (evitando el efecto convoy).

La prioridad definida externamente es asignada mediante algún mecanismo por el usuario.

El mayor problema de estos algoritmos, sin importar su definición de prioridad, es el \emph{bloqueo indefinido} o \emph{muerte por inanición}\cite{silberschatz}.
Cuando muchas tareas de alta prioridad buscan ejecutarse no permiten a aquellas de menor prioridad obtener control de la CPU.
Si esta situación se mantiene por un intervalo significativo podría afectar la experiencia del usuario y,
con un flujo estable de tareas de alta prioridad, podría evitar de forma completa la finalización de esta tarea.

Una solución a este problema es el mecanismo de \emph{envejecimiento}.
Bajo este marco se aumenta gradualmente la prioridad de las tareas que están esperando CPU, garantizando su eventual ejecución.

\section{Hilos}
El modelo de proceso discutido en clase implica que un proceso es un programa que realiza un 
solo hilo de ejecución. ``Por ejemplo, cuando un proceso está ejecutando un procesador 
de textos, un solo hilo de instrucciones se está ejecutando. Este único hilo de control permite 
que el proceso realice solo una tarea a la vez. Por ejemplo, el usuario no puede escribir simultáneamente 
caracteres y ejecutar el corrector ortográfico dentro del mismo proceso. 
La mayoría de los sistemas operativos modernos han extendido el concepto de proceso para 
permitir que un proceso tenga múltiples hilos de ejecución y así realizar más de una tarea 
al mismo tiempo.'' Esta característica es especialmente beneficiosa en sistemas multiprocesador, 
donde se pueden ejecutar múltiples hilos en paralelo.\cite{silberschatz}

``Un hilo es un flujo de control dentro de un proceso. Un proceso multihilo contiene varios 
flujos de control diferentes dentro del mismo espacio de direcciones.
Los beneficios de los mecanismos multihilo incluyen una mayor capacidad de respuesta para 
el usuario, la compartición de recursos dentro del proceso, una mayor economía y los factores 
de escalabilidad, como un uso más eficiente de múltiples núcleos de procesamiento.
Los hilos de nivel de usuario son hilos que son visibles para el programador y son desconocidos
para el kernel. El kernel del sistema operativo soporta y administra hilos de nivel de kernel. 
En general, los hilos de usuario son más rápidos de crear y gestionar que los hilos del kernel, 
porque no se requiere intervención del kernel.''\cite{silberschatz}
Existen tres tipos diferentes de modelos que relacionan los hilos de usuario y del kernel:
\begin{itemize}
	\item El modelo muchos-a-uno asigna muchos hilos de usuario a un solo hilo del kernel. 
	\item El modelo uno-a-uno asigna cada hilo de usuario a un hilo de kernel correspondiente.
	\item El modelo muchos-a-muchos multiplexa muchos hilos de usuario a un número menor o igual de hilos del kernel.
\end{itemize}

\subsection{Modelos y Ámbitos de contienda}

El modelo muchos-a-muchos presenta un esquema en que \emph{dos} planificaciones deben llevarse a cabo.
En primer lugar, los hilos de usuario creados deben ser planificados cada uno en un hilo de kernel (si quedan disponibles).
Esto se conoce como ámbito de contienda de proceso (PCS)\cite{silberschatz}, ya que quienes compiten son únicamente los hilos del proceso en cuestión.
Por otro lado, los hilos de kernel planificados sobre la(s) CPU(s) compiten con hilos de otros procesos (de todo el sistema), por lo que se conoce como ámbito de contienda de sistema (SCS).
Comúnmente, el PCS sigue un esquema de prioridades para la planificación.



\section{Memoria compartida}

Uno de los problemas en los sistemas multiproceso es controlar el acceso a la
memoria. Si todos ejecutan de manera libre y
descontrolada podrían escribir en la misma dirección, acabando eventualmente con un estado de
memoria incoherente.

Además, según como ocurra la planificación del procesador, los resultados podrían
ser diferentes cada vez, haciendo impredecible la ejecución de los procesos.

La forma en que resuelven este problema los sistemas es múltiple. Una clara
opción es dar a cada proceso un espacio de memoria determinado y prohibir el
acceso de un proceso al espacio de otro. El problema de esta solución es que
algunas veces es conveniente poder compartir memoria entre procesos, ya sea para
que se comuniquen o para dividir el trabajo.

\subsection{Problema de los lectores-escritores}
Cuando un conjunto de datos es compartido por varios procesos concurrentes, unos buscan leer y otros escribir.
Para distinguirlos llamamos \emph{lectores} a los primeros y \emph{escritores} a los segundos\cite{silberschatz}.
En este marco podemos detallar los casos que presentan una situación de riesgo y que deben ser controladas.
Si dos o más lectores acceden a los datos esto no presenta ningún inconveniente.
Sin embargo, si un escritor accede a los datos ningún otro (lector o escritor) deberá hacerlo.

\subsubsection{Variantes}
En el \emph{primer} problema de lectores-escritores,
ningún lector se mantiene en espera a menos que un escritor haya obtenido ya permiso para el objeto
compartido \cite{silberschatz}.

El \emph{segundo} problema trata del caso en que un escritor preparado realiza la operación tan pronto como sea posible \cite{silberschatz}.

Este ultimo caso es el de nuestro sistema, ya que queremos que la compra se efectué tan rápido como sea posible.

\subsection{Bloqueo y Sincronización}
Para controlar el problema se busca controlar el acceso a esta memoria.
Se busca evitar de alguna manera que dos procesos modifiquen la misma locación de manera concurrente.
Una manera es hacerlo mediante semáforos.

\subsubsection{Semáforos}
Un semáforo es una variable a la cual se accede únicamente mediante dos
operaciones atómicas: \call{wait()} y \call{signal()} \cite{silberschatz}.
Posee dos propiedades: \emph{value} que es un entero y \emph{list} una lista de bloques de control de procesos.
Una posible implementación de estas operaciones es la que se muestra en el listado \ref{lst:wait-sem}.

El segmento de código en que un proceso modifica memoria compartida se conoce como \emph{sección crítica} \cite{silberschatz}.
Mediante estas
operaciones es posible crear bloques de código de sección crítica que no tendrán
colisión con otros procesos (o eventualmente hilos) que también respeten al
semáforo.

De esta manera, si el acceso a una estructura compartida es controlado mediante
semáforos es posible asegurar la consistencia e integridad de los datos.

En el caso, el conjunto de entradas disponibles y vendidas, compartida por los
múltiples servicios de venta, podría ser visto como una estructura a sincronizar
entre ellos, asegurando que cuando uno toma una entrada ningún otro lo hace.

\lstinputlisting[float,caption=Implementaci\'on de un sem\'aforo,label={lst:wait-sem}]{../../../informe/wait.c}


\subsubsection{Interbloqueos}

Un proceso que está en espera por la liberación de un recurso\footnote{Como puede ser un semáforo} no podrá ejecutar hasta que el sistema se lo asigne.
Sin embargo, mientras se encuentra esperando puede llegar a tener otros recursos bajo su control.

Imaginemos un grafo cuyos vértices son los procesos y recursos del sistema.
Si un proceso está en espera de un recurso existirá una arista desde el proceso \(P_i\) hasta el recurso, \(R_j\).
Si un proceso posee un recurso existirá una arista desde el recurso \(R_j\) hasta este proceso, \(P_k\).

\begin{center}
\begin{tikzpicture}
\def \n {4}
\def \radius {3cm}
\def \margin {8}
  \node[draw, circle] at ({360/\n * (0)}:\radius) {$P_i$};
  \draw[->, >=latex] ({360/\n * 0+\margin}:\radius)
    arc ({360/\n * 0+\margin}:{360/\n * (1)-\margin}:\radius);
  \node[draw, circle] at ({360/\n * (1)}:\radius) {$R_j$};
  \draw[->, >=latex] ({360/\n * 1+\margin}:\radius)
    arc ({360/\n * 1+\margin}:{360/\n * (2)-\margin}:\radius);
  \node[draw, circle] at ({360/\n * (2)}:\radius) {$P_k$};
  \draw[->, >=latex] ({360/\n * 2+\margin}:\radius)
    arc ({360/\n * 2+\margin}:{360/\n * (3)-\margin}:\radius);
  \node[draw, circle] at ({360/\n * (3)}:\radius) {$R_h$};
  \draw[->, >=latex] ({360/\n * 3+\margin}:\radius)
    arc ({360/\n * 3+\margin}:{360/\n * (4)-\margin}:\radius);
\end{tikzpicture}
\end{center}

Pero si analizamos la situación en que un ciclo se dé en este grafo (con el camino \(P_i,R_j,P_k,R_h\)) encontramos un problema.
Ambos procesos esperan por un recurso, pero ninguno de los dos lo obtendrá hasta que el otro libere el que tiene en su posesión,
cosa que no hará hasta obtener el recurso que espera.
La figura A continuación mostramos un par de hilos que entran en interbloqueo esperando por dos semáforos.

\begin{center}
\begin{tabular}{|l|l|}
	Hilo A & Hilo B \\ \hline
	wait(Sem1) & wait(Sem2) \\
	wait(Sem2) & wait(Sem1) \\
	\multicolumn{2}{|c|}{// Nunca ejecuta} \\
	signal(Sem1) & signal(Sem2) \\
	signal(Sem2) & signal(Sem1) \\
	\hline
\end{tabular}
\end{center}


\chapter{Alternativas de solución}
\section{Marco general}

El sistema de ventas consiste de un servidor \emph{Web Services} el cual
proporciona una API que permite a los usuarios realizar la consulta de
entradas y la compra de las mismas.

Para realizar esta operación el sistema solicita que se esté autenticado
(creando un usuario y una contraseña) para permitir continuar con las
operaciones.

\section{Concurrencia de la transacción}
Un problema a resolver, en cuanto a lo técnico, es que la venta de
entradas funcione de manera concurrente, permitiendo a muchos clientes operar a
la vez no solo permitiendo un bajo tiempo de respuesta sino imposibilitando
fallas como vender una misma entrada a múltiples compradores.

\subsection{Bloqueo para cada operación} \label{alt-bloqueo}
Una solución posible consiste en utilizar semáforos para controlar el acceso a
una estructura principal de \emph{entradas disponibles}, única para cada evento
(ver Figura \ref{fig:bloqueocadaop}).

Cuando un usuario comience la operación elegirá una entrada a comprar, luego de
lo cual se le solicitará información de pago y demás. En el momento en que
escoge el asiento se quitará la entrada de la estructura, controlando el acceso
a ella mediante un semáforo.

\begin{figure}
\centering{
	\def\svgwidth{\linewidth}
\input{img/bloqueo.pdf_tex}
\caption{\nameref{alt-bloqueo}}
\label{fig:bloqueocadaop}
}
\end{figure}

\begin{verbatim}
GetSemaphore();
RemoveTicketIfAvailable(ticketId);
ReleaseSemaphore();
\end{verbatim}

De esta forma nos aseguramos que la estructura siempre contiene las entradas que
están disponibles y que cuando un cliente elija una entrada solo se le
permitirá continuar si ella está realmente libre.

Al mostrarle las entradas al usuario para que elija basta con consultar la lista
de entradas disponibles y pasarlas al usuario. Si en el intervalo hasta que se
decide por una la misma fue vendida la llamada a \call{RemoveTicketIfAvailable}
avisará de la situación.

\subsection{Comunicaciones indirectas - Colas de mensajes}
Otra solución posible es utilizar colas de mensajes, un tipo de estructura que provee un mecanismo de
comunicación indirecto, asincrónico, y escalable.
Distintos productos proveen este mecanismo, uno de ellos es \type{Java Message Service}.
Este tipo de software suele implementar controles transaccionales (concurrencia y recuperación)
y cuenta con mecanismos de seguridad.

Cuando un usuario comience la operación elegirá una entrada a comprar, luego de
lo cual se le solicitará información de pago y demás. Luego de seleccionar el asiento,
una estructura \type{HashMap} se encargará de gestionar la transacción, distinguiendo
asientos ocupados y libres.

Esta solución descarta el uso de semáforos porque no se comparte la estructura que maneja los mensajes.

\subsection{Estructura de entradas vendidas o pendientes} \label{concurrenthm}
Una estructura \type{HashMap} permite la inserción y la consulta en \BigO{1}. Además,
por su funcionamiento, en la mayor parte de los casos de inserción los elementos
se colocan el lugares disjuntos de memoria, y esto permite que dos hebras
inserten de forma concurrente.

\subsubsection{Un mapa concurrente}
En particular imaginamos un tipo \type{ConcurrentHashMap}, que tiene un semáforo por cada espacio del vector (Figura \ref{fig:sistema}).

Para cada inserción que se desee realizar calcula el \emph{hash} adecuado para el objeto,
y realiza la inserción en el vector solo habiendo adquirido dicho semáforo.
Devuelve el valor asociado previamente (o nulo si no existía).

Supongamos además una operación \call{putIfAbsent}.
Ella coloca un valor asociado a cierta clave \emph{solamente} si la clave no tiene ya un valor asociado.
La operación se realiza calculando el \emph{hash} y luego,
obteniendo el cerrojo para dicho espacio del vector,
revisándolo e insertando si no existe un valor asociado.
La operación tendrá como retorno el valor anterior del mapa, o nulo si no existía la clave.

\begin{verbatim}
i = CalcularIndiceConHash(clave)
semaforos[i].Esperar();
if (!mapa.contiene(clave))
  r = map.insertar(clave, valor);
else
  r = nulo;
endif
semaforos[i].Indicar();
return r
\end{verbatim}

Java presenta un primitivo que cumple con esta interfaz, que lleva el mismo nombre.

En esta alternativa se posee una estructura \type{SoldTickets}, implementada
como un \type{HashMap} concurrente.

Cuando se inserta en una estructura \type{ConcurrentHashMap} la
operación \emph{es previa a} la búsqueda \cite{javaspec}.
Esto significa que toda operación de este último tipo que retorne luego de
haber finalizado una operación de inserción operará con el último valor del mapa.

\subsubsection{Detallado de la solución}\label{detallado-chm}

\begin{figure}
\centering{
	\def\svgwidth{\linewidth}
\input{img/sistema.pdf_tex}
\caption{Se adquiere el cerrojo para escribir, pero no bloquea leer}
\label{fig:sistema}
}
\end{figure}

El mapa \type{SoldTickets} será de identificador de entrada a \type{Unit}.
Este es un tipo \emph{unitario}, es decir que existe un único posible valor de este tipo, denotado \texttt{()}.
Esto se debe a que el valor en sí es ignorado, solo importando si existe o no en el mapa algo asociado a la clave.
Se convierte así al mapa en un conjunto de valores del tipo \type{IdEntrada}.

Así, con operaciones como \call{putIfAbsent} es posible buscar una entrada y
colocarla atómicamente\cite{javaspec} en la estructura, asegurándose de que
ningún otra hebra pudo hacerlo, garantizando por lo tanto que la venta es segura.

Una ventaja de esto es que la estructura ConcurrentHashMap solo bloquea cuando
se opera con la misma clave (o al menos con un rango restringido) por lo que en
una gran mayoría de los casos es posible operar de forma paralela, en
comparación con bloquear la estructura con cada inserción, de cierta forma
haciéndolas de manera secuencial.

\subsubsection{Compra de múltiples entradas}
La enorme mayoría de los clientes realizará la compra de \(m > 1\) entradas.
Para lidiar con este caso es importante considerar que,
si no es posible la compra de las \(m\) entradas,
no queremos realizar una de \(k < m\).

Para respetar la elección del cliente sería posible bloquear durante la reserva de las entradas,
desbloqueando una vez que todas hayan sido reservadas.
El problema de esta solución es que, con una alta probabilidad, las compras concurrentes no poseen entradas
en común.

Para cada entrada solicitada se intentará su colocación en \type{SoldTickets}.
Si alguna de ellas falla deberán liberarse las reservadas anteriormente.
Esto permite que la mayor parte de las operaciones funcionen de manera concurrente.

\subsubsection{Venta a agencias}\label{venta-agencia}
Las agencias serán capaces de comprar grandes cantidades de entradas para luego vender a discreción.
Para el sistema, se comportarán igual que los compradores de múltiples entradas descritos anteriormente,
con la diferencia de que a las agencias se les exigirá un mínimo de entradas para comprar y se les hará un descuento
en su compra.

\section{Planificación de la Venta}
Al priorizar el tiempo de respuesta frente a los usuarios se procura atender a las solicitudes tan pronto como sea posible.
Sin embargo, dadas las limitaciones físicas (menos CPUs que usuarios concurrentes) es inevitable que algunas acaben encolando.

\subsection{First Come, First Serve}
Una opción posible es encolar a estas solicitudes en una cola \emph{FIFO}, como se describe en la sección \ref{fcfs}.
Esto implica que se dará trato a la solicitud mas antigua en primer lugar, por lo que el primer usuario en llegar será prioritario.
Aunque esto mejora el tiempo de respuesta entra en conflicto con criterios de mayor prioridad,
sea maximizar el ingreso total y la ocupación del estadio (ver capítulo \ref{cap:criterios}).

%\subsection{Mayor venta reciente primero}
\subsection{Mayor venta primero}\label{mayor-primero}

Con el fin de maximizar los ingresos y la cantidad de entradas vendidas se organiza la cola de solicitudes mediante prioridades.
Como prioridad de una solicitud se toma \emph{el opuesto} del valor total de la transacción.
\[ p = -c_{venta}\].
Esto garantiza que se dará servicio a la venta de mayor valor en primer lugar y, en particular,
se atenderán primero las ventas de la mayor cantidad de entradas en el caso promedio (considere entradas del mismo precio).




\section{Solución final}

La solución detallada en la sección \ref{concurrenthm} será implementada, junto con la planificación descrita en \ref{mayor-primero}.
La ventaja que ofrece esta estructura, permitiendo múltiples inserciones sin bloqueo (más allá de alguna excepción)
es grande frente a la alternativa de efectivamente hacer secuenciales a las operaciones de manipulación de la estructura.

Por otro lado, atender a las solicitudes de mayor ganancia primero permite maximizar los ingresos,
y de esta forma se resuelven los primeros 3 criterios de optimización.

Se implementa un chequeo de asientos aislados para satisfacer el cuarto criterio imposibilitando ventas que dejen un asiento en tal estado.
Finalmente, se ofrece la posibilidad de los clientes con discapacidades de autenticarse como tal ante el sistema,
permitiendo el acceso a entradas reservadas.

\subsection{SoldTickets y Event}

La solución planteada trabaja con tres estructuras de datos íntimamente relacionadas.
Por un lado tenemos el concepto de \type{Event} (listado \ref{lst:event}), que no es más que que el espectáculo descrito en la sección \ref{recurso:espectaculo}, página \pageref{recurso:espectaculo}.
Esta estructura representa no solo el nombre para el usuario (con el cuál podrá identificar al espectáculo) si no también el momento de ocurrencia.
Finalmente posee además a la estructura \type{SoldTickets} asociada.

\lstinputlisting[float,linerange={8-11,41-41}, caption=Event.java,label={lst:event}]{modelo/Event.java}

Como se describe en en la sección \ref{concurrenthm}, \type{SoldTickets} (listado \ref{lst:soldtickets}) no es mucho más que un mapa (no obstante concurrente) a un tipo unitario
\footnote{Debido a la falta de tipos unitarios en Java se elige al tipo Boolean como imagen}.
Esta estructura permite confirmar la venta de cualquier entrada y, mediante el atributo \emph{validIds} permite identificar rápidamente las entradas válidas (disponibles o no).
Nótese que, al no modificarse este segundo mapa, es posible usar un tipo no concurrente y mejorando el desempeño.

\lstinputlisting[float,linerange={9-12,112-112}, caption=SoldTickets.java,label={lst:soldtickets}]{modelo/SoldTickets.java}


\subsection{Protocolo}
Se desarrolla un pequeño protocolo para la comunicación del servidor con los sistemas que accedan a él.
En el mismo, cada respuesta al cliente se compone de lineas.
La primera es un entero que representa la cantidad de lineas adicionales que serán transmitidas.

Cada comunicación con él se realiza mediante sockets, como se indica en la figura \ref{fig:sistema}.
Con cada nueva conexión se levanta una nueva hebra de servicio a la solicitud, \type{SocketReceiver}.
Ella procesa la misma para interpretar el tipo.

\subsubsection{GET}
Estas solicitudes buscan información de algún tipo, ya sea funciones posibles o entradas disponibles para una función en particular.
\begin{center}
	\texttt{GET <id evento>}
\end{center}

En primer lugar se presentan aquellas funciones que no detalles un id de evento.
En estos casos se muestra al solicitante un listado de los eventos disponibles para la venta,
como se muestra en el algoritmo \ref{alg:listar-eventos}.

\begin{algorithm}
\caption{Listar los eventos disponibles por una solicitud GET}
\label{alg:listar-eventos}
\begin{algorithmic}[1]
	\Require \emph{eventos} es un vector de \type{Event}
\Procedure{listarEventos}{Socket s}
\State s.put(``Los eventos disponibles son'');
\ForAll{(i, evento) $\in$ eventos}
	\State s.put(``\{i\}) \{evento.getNombre()\} \{evento.getCuando()\}'');
\EndFor
\EndProcedure
\end{algorithmic}
\end{algorithm}

El algoritmo \ref{alg:listar-entradas} muestra la respuesta del sistema cuando se detalla un id de evento.
En este caso se representa, para el evento determinado, qué entras están disponibles (mediante id),
    así como su precio.
Para lograr esto se filtra de la colección de entradas válidas \emph{validIds} aquellas que ya hayan sido
    vendidas (buscándolas en \emph{asientosOcupados}).

\begin{algorithm}
\caption{Listar las entradas disponibles}
\label{alg:listar-entradas}
\begin{algorithmic}[1]
\Procedure{listar}{Socket s, Tokens solicitud}
	\State s.put(``Las entradas disponibles son: '');
	\State idEvent $\gets$ solicitud.nextToken()
	\State evento $\gets$ events[idEvento]
	\ForAll{(entrada, costo) $\in$ evento.soldTickets.entradasDisponibles()}
		\State s.put(``\{entrada\}) \$\{costo\}'');
	\EndFor
\EndProcedure
\Statex
\Require \emph{st.validIds} es un mapa de id entrada a costo (vendidas o no)
\Require \emph{st.asientosOcupados} es el mapa descrito en \ref{concurrenthm}
\Procedure{entradasDisponibles}{SoldTickets st}
	\ForAll{(idEntrada, costo) $\in$ st.validIds}
	\If{st.asientosOcupados.get(entrada) $=$ \textbf{null}}
		\State \textbf{yield} (idEntrada, costo)
	\EndIf
	\EndFor
\EndProcedure
\end{algorithmic}
\end{algorithm}

\subsubsection{POST}
Una solicitud de este tipo busca la compra de una o más entradas
Para ello debe especificar el id del evento y cada id de la entrada que busque comprar.

\begin{center}
	\texttt{POST <tipo usuario>\ <id evento>\ <id asiento1>[, <id asiento2>[, ...]]}
\end{center}

\begin{samepage}
Además especificará un tipo usuario \(t = 0, 1, 2\) con la siguiente correspondencia.
\begin{description}[align=left,noitemsep,labelindent=\parindent]
	\item[0:] Usuarios con discapacidades.
	\item[1:] Agencias.
	\item[2:] Otros usuarios.
\end{description}
\end{samepage}

Autenticarse con \(t = 1\) permite acceder a un descuento (ver \ref{venta-agencia}).
Con \(t=0\) el usuario accede a locaciones reservadas (con facilidad de acceso, etc.~).

Deberá, para el producto final, incluir adicionalmente un \emph{token} de usuario para verificar la autenticación en la plataforma.
Esto con el fin de poder dar uso a la información de pago vinculada a su cuenta.
Pero, con el fin de simplificar la simulación, se obvia este paso hasta la entrega final del producto.

Cuando el sistema recibe una solicitud POST comienza por validar los argumentos, obteniendo un \type{idEvento} y una colección de \type{idAsiento}.
Una vez obtenida la estructura \type{Event} puede obtener todas las estructuras \type{Ticket}, lo que permite calcular el valor monetario de la solictud.
Estos datos, junto al socket de comunicación, se colocan en una \type{PriorityBlockingQueue} nombrada \emph{ticketRequestQueue}.

\subsection{Seller y TicketRequest}

La aplicación posee algunos hilos de clase \type{Seller}, los cuales se encargan de procesar la cola de ventas \emph{ticketRequestQueue}.
Estos se encargan (listado \ref{lst:seller}) de tomar el próximo valor de la cola (si lo hay), bloqueándose de forma contraria, hasta que aparezca un nuevo valor.
Los datos se procesan como se describe en la sección \ref{detallado-chm}, y se muestra el método \emph{make} en el listado \ref{lst:tkrqst}.

\lstinputlisting[float,firstline=17, lastline=26,label={lst:seller},caption=Seller.java]{servidor/Seller.java}
\lstinputlisting[float,linerange={31-37,39-43,45-53,61-62},caption=TicketRequest.java,label={lst:tkrqst}]{servidor/TicketRequest.java}

\subsection{Arquitectura de hilos}\label{arq-hilos}

Para asegurar que las ventas siempre sean procesadas es necesario que,
siempre que \emph{ticketRequestQueue} no esté vacía,
exista un hilo de clase \type{Seller} en ejecución.

Por otro lado, también queremos que las nuevas conexiones al sistema sean atendidas, mediante un hilo \type{SocketReceiver}.
Para asegurar que ambas tareas puedan ser atendidas en todo momento se diseña un esquema de tres prioridades.
La misma se ve representada en la figura \ref{fig:hilos}.

Sea \(n\) la cantidad de procesadores disponibles.
Uno de los hilos del sistema es el encargado de recibir las nuevas conexiones y crear el hilo \type{SocketReceiver} correspondiente.
Siendo \(m = n - 1\) la cantidad de procesadores restante, tenemos un \emph{pool} de \(m\) hilos de prioridad media, sobre el cual se levantan estos últimos.
Finalmente, existen \(m\) hilos \type{Seller}, con el matiz de que uno de ellos tiene prioridad alta, y los demás baja\footnote{La proporción puede ser configurada}.
Lo que esto permite es que siempre exista un hilo de venta en ejecución (por su alta prioridad), procesando las solicitudes encoladas,
pero que al mismo tiempo siempre puedan ejecutar \(m - 1\) hilos \type{SocketReceiver}.

Notase que el procesamiento de un elemento de la cola de ventas es más eficiente que el procesamiento de una conexión por \type{SocketReceiver}.
Cada ítem de la cola de ventas posee una colección de entradas que deben ser colocadas en el mapa descrito en \ref{detallado-chm}.
Esta operación tiene un orden de \(n\BigO{1} = \BigO{n} \) con la venta de \(n\) entradas\footnote{El eventual bloqueo del mapa cambiará su rendimiento, pero amortizamos}.
Por otro lado, la respuesta a una solicitud por \type{SocketReceiver} implica no solo la recepción de datos por el puerto,
si no además el \emph{parseo} de la solicitud, y eventualmente el chequeo de validez de las entradas.


\begin{figure}
\centering{
	\def\svgwidth{0.6\linewidth}
\input{img/hilos.pdf_tex}
\caption{\nameref{arq-hilos}, prioridades \(\uparrow , \rightarrow , \downarrow \)}
\label{fig:hilos}
}
\end{figure}

El esquema incluso presenta ventajas en casos de ocio.
En los casos en que la cola de ventas sea vacía, \(m\) hilos procesarán las conexiones activas.
Cuando no existan conexiones activas, \(m\) hilos procesarán rápidamente las ventas encoladas.
Finalmente, cuando el sistema no posea conexiones ni ventas encoladas, no existirán hilos \type{SocketReceiver} y los hilos \type{Seller} se encontrarán bloqueados.
En tal caso el sistema solo tendrá un hilo de ejecución (\emph{main thread}) encargado de atender las nuevas conexiones que pudieran llegar.



\chapter{Escenarios de Prueba}
\section{Compradores reales}

Para simular la interacción con usuarios reales se genera un archivo de pruebas que en tiempos aleatorios
presenta un nuevo pedido por parte de un cliente.
La cantidad de entradas es también generada de forma aleatoria, entre 1 y 10.
El tipo de usuarios en este caso es \(t=2\), y no hay asientos reservados para clientes con discapacidades.

\subsubsection{Tiempo de respuesta}
El tiempo entre la solicitud de una entrada y la respuesta correspondiente se conoce como el \emph{tiempo de respuesta}.
Para cada solicitud realizada por los clientes se medirá el tiempo de respuesta.
Finalizada la simulación se calculará el promedio de este.

Además del promedio se medirá la relación entre \(t_i\) y \(n_i\), donde \(t_i\) es el tiempo de respuesta para el i-ésimo
cliente y \(n_i\) es la cantidad de entradas que solicita.

Se distinguirá para este último caso entre las ventas con éxito (considerando solo las interacciones que finalicen en concretar la venta)
y las demás interacciones (que finalizan con un error por \emph{entrada no disponible}).

\section{Compradores concurrentes aleatorios} \label{escenario-tiempoespera}
Considerando un único espectáculo, se introducen clientes al sistema.
Cada uno de ellos tiene una cantidad de entradas que desea comprar
(sea \(n_i\)), sin importar su ubicación.

Estos intentan realizar la compra de asientos cualesquiera en el espectáculo.
En el caso de que la compra no sea exitosa (otro cliente adquirió primero) se procederá a intentarlo otra vez.
Se introducirán en total \(m\) clientes, donde \(N = \sum_{i=1}^{m} n_i\)
es la cantidad total de entradas disponibles para el evento\footnote{
	Para cada cliente se calcula un número aleatorio con distribución uniforme.
	De esta forma 
	\(n_i \in [1, 8] \wedge m \in [\frac{N}{8}, N]\).
}.
De esta forma aseguramos que cada cliente logra finalizar su interacción habiendo adquirido las entradas que deseaba.

\begin{figure}
\end{figure}
\begin{figure}
\centering{
	\def\svgwidth{\linewidth}
\input{img/escenario1.pdf_tex}
\caption{\nameref{escenario-tiempoespera}}
\label{fig:escenario1}
}
\end{figure}

En esta interacción los clientes ingresan de forma concurrente al sistema, y su orden particular dependerá de la carga
particular del sistema en que ejecute la planificación.

\subsubsection{Tiempo de respuesta}
Una vez más se mide el tiempo de respuesta de los usuarios, realizando también un análisis de la relación entre el tiempo
de respuesta y la cantidad de entradas solicitada.

La hipótesis es que las ventas de mayor valor tendrán un tiempo de respuesta relativamente menor.
A pesar de su mayor tiempo de procesamiento serán atendidas previamente.


%%%%%%%%%%%%%%%%%%%%%%%%%%%%%%%%%%%%%%%%%%%%%%%%%%%%%%%%%%%%%%%%%%%%%%%%%%%%%%%%
\section{Carga del sistema}
Para un evento con \(n\) entradas en total se introducen \(m\) clientes al sistema, donde \(n < m\).
Al igual que en la sección \ref{escenario-tiempoespera}, los clientes intentarán
la compra de un conjunto aleatorio de entradas, re-intentando en caso de una falla.

\subsubsection{Tiempo de respuesta}
Al igual que en la sección \ref{escenario-tiempoespera}, medimos el tiempo de respuesta medio de las solicitudes,
y su relación para cada cliente con la cantidad de entradas.


\chapter{Conclusiones}
\section{\nameref{escenario-tiempoespera}}

Como primera, y obvia, conclusión, se observa que, efectivamente, se realiza una venta de \(N\) entradas, y no más ni menos,
por lo que no se observan errores de lógica.

En el caso de la simulación ejecutada, tenemos \(N = 10000\), y la cantidad de clientes que compran cada número de entradas
se presenta en la tabla \ref{table:3}.
\begin{table}[H]
\centering
\begin{tabular}{ccc}
	N\textsuperscript{o} de entradas & N\textsuperscript{o} de clientes & Tiempo medio de respuesta (ms)  \\
\hline
1 &  255 & 16.434088\\
2 &  295 & 18.517253 \\
3 &  261 & 13.974438\\
4 &  271 & 17.137831\\
5 &  270 & 15.220248\\
6 &  263 & 17.686285\\
7 &  296 & 16.506400 \\
8 &  286 & 16.822844\\
\end{tabular}
\caption{Desglose por cantidad de entradas}
\label{table:3}
\end{table}


Los tiempos de respuesta fueron medidos en cada interacción para luego calcular su promedio, según algunos criterios.
En primer caso se analiza el tiempo de respuesta promedio para toda transacción,
incluyendo las comunicaciones con el sistema de tipo ``consulta de entradas'', así como las ventas fallidas.
Por otro lado se analiza el tiempo de respuesta de las transacciones que acaban con una venta concretada.
Los resultados se presentan en la tabla \ref{table:2}.
Además se presenta en la tabla \ref{table:3} el tiempo medio para cada cantidad de entradas.

\begin{table}[H]
\centering
\begin{tabular}{l|r}
	{} &  Tiempo de respuesta medio (ms) \\
	\hline
Normal & 2.747720 \\
Transacciones exitosas  & 25.564688 \\
\end{tabular}
\caption{Tiempo de respuesta medio}
\label{table:2}
\end{table}

De estos datos se concluye que el tiempo medio de respuesta es independiente de la cantidad de entradas adquiridas.
No es posible concluir que a mayor cantidad de entradas el tiempo de respuesta disminuya, pero observamos que el tiempo tampoco es proporcional a la cantidad de entradas.
Es probable que el mayor tiempo de procesamiento sea contrarrestado por el hecho de darle atención temprana a las compras de mayor valor económico.

Estos resultados se presentan en la figura \ref{fig:respuesta_concurrente}, mostrando en rojo el valor medio.

\begin{figure}[H]
\centering{
        \includegraphics[width=0.7\textwidth]{grafica_1}
\caption{Tiempo de respuesta vs. cantidad de entradas}
}
\label{fig:respuesta_concurrente}
\end{figure}

\section{Carga del sistema}

En el caso de la simulación ejecutada, tenemos \(N = 10000\), y la misma cantidad de clientes. La cantidad de clientes que compran según
las distintas cantidades de entradas que se presentan en la tabla \ref{table:carga}.

\begin{table}[H]
\centering
\begin{tabular}{ccc}
	N\textsuperscript{o} de entradas & N\textsuperscript{o} de clientes & Tiempo medio de respuesta (ms)  \\
\hline
1                 & 1206 &            24.575421 \\
2                 & 1259 &            21.579129 \\
3                 & 1253 &            26.233702 \\
4                 & 1266 &            24.669094 \\
5                 & 1260 &            25.211271 \\
6                 & 1274 &            27.186334 \\
7                 & 1254 &            23.700102 \\
8                 & 1228 &            25.710476 \\
\end{tabular}
\caption{Desglose por cantidad de entradas}
\label{table:carga}
\end{table}

Una vez más observamos que el tiempo de respuesta es independiente de la cantidad de entradas solicitadas.
Los resultados se presentan en la tabla \ref{table:carga2}.
Además se presenta en la tabla \ref{table:carga} el tiempo medio para cada cantidad de entradas.
\begin{table}[H]
\centering
\begin{tabular}{l|r}
	{} &  Tiempo de respuesta medio (ms) \\
	\hline
Normal & 7.101191 \\
Transacciones exitosas  & 24.801971 \\
\end{tabular}
\caption{Tiempo de respuesta medio}
\label{table:carga2}
\end{table}

Al analizar los tiempos de respuesta se puede ver que en comparación con la simulación previa, estos no varían al aumentar significativamente la cantidad de clientes.
Estos resultados se presentan en la figura \ref{fig:respuesta_carga}, mostrando en rojo el valor medio.

\begin{figure}[H]
\centering{
        \includegraphics[width=0.7\textwidth]{graficaMuchos1}
\caption{Tiempo de respuesta vs. cantidad de entradas}
}
\label{fig:respuesta_carga}
\end{figure}

\section{Compradores reales}

En el caso de la simulación ejecutada, tenemos \(N = 5000\), y \(M = 10000\) clientes. La cantidad de clientes que compran según
las distintas cantidades de entradas, y el tiempo medio para cada cantidad de entradas, se presentan en la tabla \ref{table:carga}.
Los compradores que realizan consultas de qué entradas se encuentran disponibles se muestran al final de la misma.

\begin{table}[H]
\centering
\begin{tabular}{ccc}
	N\textsuperscript{o} de entradas & N\textsuperscript{o} de clientes & Tiempo medio de respuesta (ms)  \\
\hline
1                 & 661 &            0.396308 \\
2                 & 640 &            0.824615 \\
3                 & 618 &            2.713671 \\
4                 & 685 &            -- \\
5                 & 640 &            0.364268 \\
6                 & 690 &            0.556214 \\
7                 & 624 &            -- \\
8                 & 730 &            -- \\
9                 & 649 &            2.453426 \\
10                & 684 &            -- \\
Consulta                 & 3380 &           1.338888 \\
\end{tabular}
\caption{Desglose por cantidad de entradas}
\label{table:reales}
\end{table}

Se dio en esta simulación que ninguno de los clientes que intentó comprar 4, 7, 8 o 10 entradas tuvo éxito en su reserva,
por lo tanto no se pudo registrar un tiempo de respuesta para estos clientes.
Se puede observar que no hay una dependencia entre el tiempo de respuesta y la cantidad de entradas solicitadas, lo cual
nos lleva a concluir que son efectivamente variables independientes, y que las fluctuaciones en el tiempo de respuesta se deben
al azar en los tiempos de llegada de cada cliente y la cantidad de clientes concurrentes en cada momento.

Los resultados se presentan en la tabla \ref{table:carga2} y en la figura, mostrando en rojo el valor medio \ref{fig:respuesta_reales},

\begin{table}[H]
\centering
\begin{tabular}{l|r}
	{} &  Tiempo de respuesta medio (ms) \\
	\hline
Normal & 0.649355 \\
Transacciones exitosas  & 1.331675 \\
\end{tabular}
\caption{Tiempo de respuesta medio}
\label{table:reales2}
\end{table}

\begin{figure}[H]
\centering{
        \includegraphics[width=0.7\textwidth]{grafica_reales}
\caption{Tiempo de respuesta vs. cantidad de entradas}
}
\label{fig:respuesta_reales}
\end{figure}

Una observación interesante a partir de la gráfica es que a pesar de que el tiempo de respuesta medio para las solicitudes de consulta es bajo
y similar a los tiempos de las de reserva, muchas consultas tuvieron un tiempo de respuesta mucho mayor. Esto refleja cómo el sistema
le da mayor prioridad a las solicitudes de reserva (ya que son las que generan ingreso) que a las de consulta, y por lo tanto demora las respuestas
de consulta con tal de responder rápido a las reservas.

\section{Conclusiones generales}

El uso de la estructuras de datos concurrente \type{ConcurrentHashMap} proporcionó una abstracción que permitió el desarrollo eficaz
sin mayor preocupación por el control del acceso a la estructura compartida.
Además, a diferencia de un cerrojo mutex para proteger una estructura común, nuestra estructura permite el acceso concurrente en la mayor parte de los casos,
aspecto que se observa en los tiempos de espera medidos.
La eficiencia de dar uso a la ejecución en paralelo mejoró el desempeño del sistema significativamente.

Para la instalación final del producto deberán ser implementada algunas características menores que se obviaron para presentar las simulaciones más relevantes.
En primer lugar los datos del programa (eventos, esquema de distribución de asientos del estadio, descuentos y proporción de hilos de venta contra solicitud) serán leídos
de datos locales, permitiendo que sean sobrescritos fácilmente.
Además se implementará el sistema de autenticación de usuarios, así como la opción de pago online mediante uso de tarjetas o puestos de venta.
Con un usuario autenticado será posible además poseer datos del mismo, como una dirección de correo mediante la cual comunicar sobre sus compras,
enviándole entradas digitales.


\appendix

\chapter{Diagrama de clases}

\begin{figure}[H]
\centering{
      \makebox[\textwidth][c]{\includegraphics[width=1.1\textwidth]{DiagramaClasesXD}}
    \caption{Diagrama de clases}
}
\end{figure}

\chapter{Máquinas virtuales}

\section{Hardware pass-through}
El método de virtualización de PCI passthrough sigue la línea que marcan la virtualización asistida por hardware (Intel VT-x y AMD-V) y la para-virtualización, con el fin común de permitir mayor performance a la máquina virtual.
En este caso, se da al sistema huésped acceso dedicado y exclusivo al dispositivo real del anfitrión, permitiéndole comunicarse directamente con él. La ventaja es clara, y es que de esta manera el costo (overhead) del acceso al hardware es mínimo.
Las desventajas de este sistema, sin embargo, son varias, por lo que debe analizarse la situación puntual para decidir si es adecuada. En primer lugar, se debe poseer soporte del hardware (Intel VT-d y AMD-Vi). Además, según como estén agrupados los dispositivos (según sus conexiones a qué controladoras), podemos perder acceso a varios de ellos ya que es necesario dar al sistema huésped el control total de un grupo, y no es posible elegir dispositivos individuales.

\section{Configuración de las MVs}
Se crean máquinas que corren Ubuntu-64bits, con \texttt{openjdk-11.0.2}, para permitir la ejecución de archivos jar. Se adjunta la imagen de la misma, y estos archivos.
Luego se configura la red para permitir la comunicación entre sí. Para esto se crea una red Host-only (\texttt{vboxnet0}), y se agrega a cada máquina un adaptador conectado a ella. Con \texttt{ip addr show} encontramos la ip de cada máquina, para abrir el puerto adecuado. Esto se realiza corriendo los programas java servidor.jar en una de ellas y \texttt{cliente.jar <ip>} en las demás.



\begin{thebibliography}{1}
\raggedright
\bibitem{silberschatz}
Silberschatz, A., Galvin, P. B., Gagne, G., \& Allende, J. S.
(2006).
\textit{Fundamentos de sistemas operativos}.
Madrid: McGraw-Hill.

\bibitem{catedra-planificacion}
Montañes, J., Barrios, E.
(2015).
\textit{Planificación de Procesos}.
Montevideo: FIT.

\bibitem{javaspec}
Java™ Platform, Standard Edition 8 API Specification.
(2019).
Disponible en: \texttt{https://docs.oracle.com/javase/8/docs/api/index.html}
\end{thebibliography}

\end{document}
